\documentclass{article}

% Language setting
% Replace `english' with e.g. `spanish' to change the document language
\usepackage[english]{babel}

% Set page size and margins
% Replace `letterpaper' with `a4paper' for UK/EU standard size
\usepackage[a4paper,top=2cm,bottom=2cm,left=3cm,right=3cm,marginparwidth=1.75cm]{geometry}

\begin{document}
\hfill Jan Romanovský
\begin{center}
    \textbf{Discoland}
\end{center}


El miércoles pasado vi la producción teatral Discoland del teatro Na zábradlí. La producción de Petr Erbes y Boris Jedinák guía al espectador tras una serie de escenas, donde un grupo de guionistas intenta conseguir material para su juego teatral sobre el club Discoland y Ivan Jonák. El espectador ve el testimonio de varias personas con varias relaciones con Discoland, por ejemplo el cocinero, el director o la bailarina en barra. Algunos se acuerdan de la época como el mejor tiempo de su vida, otros describen a Ivan Jonák como un ladrón tonto, que se comportó como en el oeste salvaje.

El juego tuvo un mal comienzo, cuando apareció una testigo sin ninguna introducción – así que entendí la esencia del juego después de unos 20 minutos. Tampoco ayudó que la primera testigo saliera de su papel y yo no sabía si era una intención dramática o la actriz simplemente no podía dejar de reírse de sus propias réplicas. Los otros actores estuvieron buenos. Destacaron Jiří Vyorálek, con su explosivo discurso sobre su descontento con la sociedad checa después de la revolución de 1989 y Václav Vašák, con su baile acrobático en barra, durante el cual habló sin ninguna dificultad respiratoria.

Lo que también me gustó fueron el vestuario y la escenografía. La escenografía era minimalista, que no molestaba el juego nada. El vestuario era fiel a la realidad, como por ejemplo el chándal rojo de Jiří Vyorálek, que me puedo imaginar en un jubilado portero de discoteca.

En general la producción me gustó, pero diría que no es para todos. Sin duda lo recomiendo a todos que saben algo de Ivan Jonák y tienen una noche libre y recomendaría leer el programa antes del juego.
\end{document}
